%==============================================================================
% Chapitre 1 - Référentiels
%==============================================================================

\section{Pourquoi les référentiels sont importants}

Lorsqu'un observateur décrit un mouvement, il le fait toujours par rapport à
quelque chose.

Ce « quelque chose » est appelé un \textbf{référentiel}.

Sans référentiel, il est impossible de définir correctement :

\begin{itemize}
    \item une position ;
    \item une vitesse ;
    \item une accélération ;
    \item une trajectoire.
\end{itemize}

Par exemple, une personne assise dans un train peut affirmer être immobile.
Pourtant, une personne observant ce même train depuis le quai indiquera que
cette personne se déplace à plusieurs dizaines de kilomètres par heure.

Les deux observations sont correctes.

Elles sont simplement effectuées dans des référentiels différents.

En balistique, le choix du référentiel est fondamental car il conditionne la
façon dont les mouvements sont observés et calculés.

\section{Intuition}

Une erreur fréquente consiste à considérer qu'un objet possède une vitesse
absolue.

En réalité, toute vitesse est relative à un référentiel.

Un passager assis dans un train roulant à 100~km/h possède :

\begin{itemize}
    \item une vitesse nulle par rapport au train ;
    \item une vitesse de 100~km/h par rapport au sol.
\end{itemize}

Il n'existe aucune contradiction.

Les deux mesures décrivent le même phénomène depuis deux points de vue
différents.

\section{Exemple : tir entre deux véhicules}

Considérons deux véhicules :

\begin{itemize}
    \item roulant à 10~km/h ;
    \item dans la même direction ;
    \item à vitesse constante ;
    \item séparés latéralement de 50~m.
\end{itemize}

Un tireur situé sur le premier véhicule vise un point fixe du second véhicule.

Intuitivement, on pourrait penser qu'il faut viser en avant de la cible puisque
celle-ci se déplace.

Pourtant, ce n'est pas le cas.

Au moment du tir :

\begin{itemize}
    \item le projectile hérite de la vitesse du véhicule tireur ;
    \item le véhicule cible possède la même vitesse.
\end{itemize}

Dans le référentiel des véhicules :

\begin{itemize}
    \item le tireur est immobile ;
    \item la cible est immobile ;
    \item la distance reste constante.
\end{itemize}

Le problème devient identique à un tir entre deux objets fixes.

Aucune correction de visée n'est nécessaire.

Cet exemple illustre l'importance du choix du référentiel.

\section{Exemple : une balle lancée dans un train}

Une expérience classique consiste à lancer une balle verticalement dans un
train roulant à vitesse constante.

Pour un observateur situé dans le train :

\begin{itemize}
    \item la balle monte ;
    \item la balle redescend ;
    \item elle retombe dans la main du lanceur.
\end{itemize}

Pour un observateur situé à l'extérieur :

\begin{itemize}
    \item la balle avance avec le train ;
    \item elle décrit une trajectoire courbe.
\end{itemize}

Les deux observations sont correctes.

La différence provient uniquement du référentiel utilisé.

\section{Référentiel terrestre}

Pour la plupart des applications balistiques, le référentiel utilisé est celui
lié à la surface de la Terre.

Les positions sont exprimées par rapport au sol.

Les vitesses sont exprimées par rapport au terrain.

Les cibles sont généralement considérées fixes dans ce référentiel.

Pour les tirs de courte et moyenne portée, cette approximation est largement
suffisante.

\section{Limites du référentiel terrestre}

La Terre tourne sur elle-même.

Par conséquent, le référentiel terrestre n'est pas rigoureusement inertiel.

Cette rotation engendre plusieurs phénomènes :

\begin{itemize}
    \item l'effet Coriolis ;
    \item l'effet Eötvös ;
    \item la courbure terrestre ;
    \item les variations du champ gravitationnel.
\end{itemize}

Pour les tirs courants, ces effets sont faibles.

Pour les tirs à longue distance ou pour l'artillerie, ils deviennent
mesurables.

\section{Référentiel galiléen}

Un référentiel galiléen est un référentiel dans lequel les lois de Newton
s'appliquent directement.

Dans un tel référentiel :

\begin{itemize}
    \item un objet sans force appliquée conserve sa vitesse ;
    \item une accélération est nécessaire pour modifier son mouvement.
\end{itemize}

En pratique, aucun référentiel réel n'est parfaitement galiléen.

Cependant, pour de nombreuses applications, certains référentiels peuvent être
considérés comme suffisamment proches d'un référentiel galiléen.

\section{Référentiel tournant}

Lorsqu'un référentiel est en rotation, des forces dites \emph{apparentes}
apparaissent.

Ces forces ne résultent pas d'une interaction physique directe.

Elles proviennent du choix du référentiel.

Les principales sont :

\begin{itemize}
    \item la force centrifuge ;
    \item la force de Coriolis.
\end{itemize}

Ces forces sont indispensables pour décrire correctement les mouvements
observés depuis la surface de la Terre.

\section{Application à la balistique}

Les référentiels interviennent dans :

\begin{itemize}
    \item le calcul des trajectoires ;
    \item la modélisation des cibles mobiles ;
    \item les tirs embarqués ;
    \item les tirs depuis des aéronefs ;
    \item les tirs longue distance ;
    \item les simulateurs.
\end{itemize}

La compréhension des référentiels constitue donc l'une des bases de toute
étude balistique.

\section{À retenir}

\begin{aretenir}

\begin{itemize}
    \item Toute vitesse est relative à un référentiel.
    \item Un même mouvement peut être observé différemment selon le référentiel choisi.
    \item Le référentiel terrestre est généralement utilisé en balistique.
    \item La rotation de la Terre introduit des effets supplémentaires.
    \item Les référentiels sont indispensables pour comprendre les trajectoires, les cibles mobiles et les phénomènes tels que l'effet Coriolis.
\end{itemize}

\end{aretenir}
